\documentclass[12pt,a4paper]{ctexart}

% ============================================================
%                       宏包设置
% ============================================================

\usepackage{amsmath, amssymb, amsfonts}
\usepackage{geometry}
\usepackage{fancyhdr}
\usepackage{graphicx}
\usepackage{xcolor}
\usepackage{setspace}
\usepackage[most]{tcolorbox}
\usepackage{enumitem}

% ============================================================
%                 清新蓝绿色配色
% ============================================================

% 主色：清爽青绿色
\definecolor{SZURed}{HTML}{2F9E8F}

% 题目区域：浅薄荷绿
\definecolor{SZULightRed}{HTML}{EAF7F4}

% 答案区域：浅天空蓝
\definecolor{SZULightRose}{HTML}{EEF7FC}

% 边框：柔和蓝绿色
\definecolor{SZUSoftBorder}{HTML}{9CCFC8}

% 强调文字：沉稳蓝绿色
\definecolor{SZUTextRed}{HTML}{256B72}

% ============================================================
%                       页面设置
% ============================================================

\geometry{
    left=2.5cm,
    right=2.5cm,
    top=2.5cm,
    bottom=2.5cm
}

\setlength{\parindent}{0pt}
\setlength{\parskip}{0.6em}
\linespread{1.15}

% ============================================================
%                       课程信息
% ============================================================

\newcommand{\coursename}{概率论与数理统计}
\newcommand{\department}{深圳大学生物医学工程}
\newcommand{\homeworknumber}{作业 1}

% ============================================================
%                  学生信息 —— 请学生修改
% ============================================================

\newcommand{\studentname}{请填写姓名}
\newcommand{\studentid}{请填写学号}

% ============================================================
%                       页眉页脚
% ============================================================

\pagestyle{fancy}
\fancyhf{}

\lhead{\textcolor{SZUTextRed}{\coursename}}
\rhead{\textcolor{SZUTextRed}{\homeworknumber}}
\cfoot{\textcolor{SZUTextRed}{\thepage}}

\renewcommand{\headrulewidth}{0.4pt}
\renewcommand{\headrule}{
    \hbox to\headwidth{
        \color{SZUSoftBorder}\leaders\hrule height \headrulewidth\hfill
    }
}

\setlength{\headheight}{15pt}

% ============================================================
%                    题目编号与环境设置
% ============================================================

\newcounter{problem}

% -------------------------
% 题目环境
% -------------------------

\newenvironment{problem}[1][]
{
    \stepcounter{problem}
    \begin{tcolorbox}[
        title={题目 \theproblem\if&#1&\relax\else\quad[#1]\fi},
        colback=SZULightRed,
        colframe=SZURed,
        coltitle=white,
        fonttitle=\bfseries,
        boxrule=0.8pt,
        arc=2mm,
        left=3mm,
        right=3mm,
        top=2.5mm,
        bottom=2.5mm,
        breakable
    ]
}
{
    \end{tcolorbox}
}

% -------------------------
% 答案环境
% -------------------------

\newenvironment{answer}
{
    \begin{tcolorbox}[
        title={学生答案},
        colback=SZULightRose,
        colframe=SZUSoftBorder,
        coltitle=SZUTextRed,
        fonttitle=\bfseries,
        boxrule=0.7pt,
        arc=2mm,
        left=3mm,
        right=3mm,
        top=2.5mm,
        bottom=2.5mm,
        breakable
    ]
}
{
    \end{tcolorbox}
}

% ============================================================
%                         正文
% ============================================================

\begin{document}

% ============================================================
%                         标题
% ============================================================

\begin{center}

    {\LARGE\bfseries\textcolor{SZURed}{\coursename}}

    \vspace{0.4em}

    {\large \department}

    \vspace{1em}

    {\Large\bfseries\textcolor{SZUTextRed}{\homeworknumber}}

\end{center}

\vspace{1em}

% ============================================================
%                       学生信息
% ============================================================

\begin{center}

\begin{tabular}{rl@{\hspace{3cm}}rl}

\textcolor{SZUTextRed}{\bfseries 姓名：}
&
\studentname
&
\textcolor{SZUTextRed}{\bfseries 学号：}
&
\studentid

\end{tabular}

\end{center}

\vspace{0.7em}

{\color{SZUSoftBorder}\hrule}

\vspace{1.5em}



% ============================================================
% ============================================================
%
%                           题目 1
%
% ============================================================
% ============================================================

\begin{problem}[课本习题2]

设 $A,B,C$ 为三个事件，用 $A,B,C$ 的运算关系表示下列各事件：

\begin{enumerate}[label=(\arabic*)]
    \item $A$ 发生，$B$ 与 $C$ 不发生。
    \item $A$ 与 $B$ 都发生，而 $C$ 不发生。
    \item $A,B,C$ 中至少有一个发生。
    \item $A,B,C$ 都发生。
    \item $A,B,C$ 都不发生。
    \item $A,B,C$ 中不多于一个发生。
    \item $A,B,C$ 中不多于两个发生。
    \item $A,B,C$ 中至少有两个发生。
\end{enumerate}

\end{problem}


% ============================================================
%
%              ↓↓↓↓↓↓ 学生答案区域 ↓↓↓↓↓↓
%
%          请只在 answer 环境中填写你的答案
%
% ============================================================

\begin{answer}

% ==========================================================
%                 请从下一行开始填写答案
% ==========================================================






% ==========================================================
%                 请在上一行结束填写答案
% ==========================================================

\end{answer}

% ============================================================
%
%              ↑↑↑↑↑↑ 学生答案区域 ↑↑↑↑↑↑
%
% ============================================================


\vspace{1em}



% ============================================================
% ============================================================
%
%                           题目 2
%
% ============================================================
% ============================================================

\begin{problem}[课本习题3(2)]

已知 $P(A)=\dfrac12$，$P(B)=\dfrac13$，$P(C)=\dfrac15$，$P(AB)=\dfrac1{10}$，$P(AC)=\dfrac1{15}$，$P(BC)=\dfrac1{20}$，$P(ABC)=\dfrac1{30}$，求 $A\cup B,\bar{A}\bar{B},A\cup B\cup C,\bar{A}\bar{B}\bar{C},\bar{A}\bar{B}C,\bar{A}\bar{B}\cup C$ 的概率。

\end{problem}


% ============================================================
%
%              ↓↓↓↓↓↓ 学生答案区域 ↓↓↓↓↓↓
%
% ============================================================

\begin{answer}

% ==========================================================
%                 请从下一行开始填写答案
% ==========================================================





% ==========================================================
%                 请在上一行结束填写答案
% ==========================================================

\end{answer}

% ============================================================
%
%              ↑↑↑↑↑↑ 学生答案区域 ↑↑↑↑↑↑
%
% ============================================================


\vspace{1em}



% ============================================================
% ============================================================
%
%                           题目 3
%
% ============================================================
% ============================================================

\begin{problem}[课本习题6]

在房间里有10个人，分别佩戴从1号到10号的纪念章，任选3人记录其纪念章的号码。
\begin{enumerate}[label=(\arabic*)]
    \item 求最小号码为5的概率。
    \item 求最大号码为5的概率。
\end{enumerate}

\end{problem}


% ============================================================
%
%              ↓↓↓↓↓↓ 学生答案区域 ↓↓↓↓↓↓
%
% ============================================================

\begin{answer}

% ==========================================================
%                 请从下一行开始填写答案
% ==========================================================





% ==========================================================
%                 请在上一行结束填写答案
% ==========================================================

\end{answer}

% ============================================================
%
%              ↑↑↑↑↑↑ 学生答案区域 ↑↑↑↑↑↑
%
% ============================================================


\vspace{1em}



% ============================================================
% ============================================================
%
%                           题目 4
%
% ============================================================
% ============================================================

\begin{problem}[课本习题10]

在11张卡片上分别写上 \textit{probability} 这11个字母，从中任意连抽7张，求其排列结果为 \textit{ability} 的概率。

\end{problem}


% ============================================================
%
%              ↓↓↓↓↓↓ 学生答案区域 ↓↓↓↓↓↓
%
% ============================================================

\begin{answer}

% ==========================================================
%                 请从下一行开始填写答案
% ==========================================================





% ==========================================================
%                 请在上一行结束填写答案
% ==========================================================

\end{answer}

% ============================================================
%
%              ↑↑↑↑↑↑ 学生答案区域 ↑↑↑↑↑↑
%
% ============================================================


\vspace{1em}



% ============================================================
% ============================================================
%
%                           题目 5
%
% ============================================================
% ============================================================

\begin{problem}[课本习题12]

50只铆钉随机地取来用在10个部件上，其中有3只铆钉强度太弱。每个部件用3只铆钉。若将3只强度太弱的铆钉都装在一个部件上，则这个部件强度就太弱。问发生一个部件强度太弱的概率是多少？

\end{problem}


% ============================================================
%
%              ↓↓↓↓↓↓ 学生答案区域 ↓↓↓↓↓↓
%
% ============================================================

\begin{answer}

% ==========================================================
%                 请从下一行开始填写答案
% ==========================================================





% ==========================================================
%                 请在上一行结束填写答案
% ==========================================================

\end{answer}

% ============================================================
%
%              ↑↑↑↑↑↑ 学生答案区域 ↑↑↑↑↑↑
%
% ============================================================


\vspace{1em}



% ============================================================
% ============================================================
%
%                           题目 6
%
% ============================================================
% ============================================================

\begin{problem}[课本习题14(1)]

已知 $P(\bar{A})=0.3$，$P(B)=0.4$，$P(A\bar{B})=0.5$，求条件概率 $P(B\mid A\cup \bar{B})$。

\end{problem}


% ============================================================
%
%              ↓↓↓↓↓↓ 学生答案区域 ↓↓↓↓↓↓
%
% ============================================================

\begin{answer}

% ==========================================================
%                 请从下一行开始填写答案
% ==========================================================





% ==========================================================
%                 请在上一行结束填写答案
% ==========================================================

\end{answer}

% ============================================================
%
%              ↑↑↑↑↑↑ 学生答案区域 ↑↑↑↑↑↑
%
% ============================================================


\vspace{1em}



% ============================================================
% ============================================================
%
%                           题目 7
%
% ============================================================
% ============================================================

\begin{problem}[课本习题18]

某人忘记了电话号码的最后一个数字，因而他随意地拨号。求他拨号不超过三次而接通所需电话的概率。若已知最后一个数字是奇数，那么此概率是多少？

\end{problem}


% ============================================================
%
%              ↓↓↓↓↓↓ 学生答案区域 ↓↓↓↓↓↓
%
% ============================================================

\begin{answer}

% ==========================================================
%                 请从下一行开始填写答案
% ==========================================================





% ==========================================================
%                 请在上一行结束填写答案
% ==========================================================

\end{answer}

% ============================================================
%
%              ↑↑↑↑↑↑ 学生答案区域 ↑↑↑↑↑↑
%
% ============================================================


\vspace{1em}



% ============================================================
% ============================================================
%
%                           题目 8
%
% ============================================================
% ============================================================

\begin{problem}[课本习题22]

一学生接连参加同一课程的两次考试：第一次及格的概率为 $p$，若第一次及格则第二次及格的概率也为 $p$；若第一次不及格则第二次及格的概率为 $p/2$。
\begin{enumerate}[label=(\arabic*)]
    \item 若至少有一次及格则他能取得某种资格，求他取得该资格的概率。
    \item 若已知他第二次已经及格，求他第一次及格的概率。
\end{enumerate}

\end{problem}


% ============================================================
%
%              ↓↓↓↓↓↓ 学生答案区域 ↓↓↓↓↓↓
%
% ============================================================

\begin{answer}

% ==========================================================
%                 请从下一行开始填写答案
% ==========================================================





% ==========================================================
%                 请在上一行结束填写答案
% ==========================================================

\end{answer}

% ============================================================
%
%              ↑↑↑↑↑↑ 学生答案区域 ↑↑↑↑↑↑
%
% ============================================================


\vspace{1em}



% ============================================================
% ============================================================
%
%                           题目 9
%
% ============================================================
% ============================================================

\begin{problem}[课本习题23]

将两信息分别编码为 $A$ 和 $B$ 传送出去，接收站收到时，$A$ 被误收作 $B$ 的概率为 $0.02$，而 $B$ 被误收作 $A$ 的概率为 $0.01$。信息 $A$ 与信息 $B$ 传送的频繁程度为 $2:1$。若接收站收到的信息是 $A$，问原发信息是 $A$ 的概率是多少？

\end{problem}


% ============================================================
%
%              ↓↓↓↓↓↓ 学生答案区域 ↓↓↓↓↓↓
%
% ============================================================

\begin{answer}

% ==========================================================
%                 请从下一行开始填写答案
% ==========================================================





% ==========================================================
%                 请在上一行结束填写答案
% ==========================================================

\end{answer}

% ============================================================
%
%              ↑↑↑↑↑↑ 学生答案区域 ↑↑↑↑↑↑
%
% ============================================================


\vspace{1em}



% ============================================================
% ============================================================
%
%                           题目 10
%
% ============================================================
% ============================================================

\begin{problem}[课本习题26]

病树的主人外出，委托邻居浇水，设已知如果不浇水，树死去的概率为 $0.8$。若浇水则树死去的概率为 $0.15$。有 $0.9$ 的把握确定邻居会记得浇水。
\begin{enumerate}[label=(\arabic*)]
    \item 求主人回来树还活着的概率。
    \item 若主人回来树已死去，求邻居忘记浇水的概率。
\end{enumerate}

\end{problem}


% ============================================================
%
%              ↓↓↓↓↓↓ 学生答案区域 ↓↓↓↓↓↓
%
% ============================================================

\begin{answer}

% ==========================================================
%                 请从下一行开始填写答案
% ==========================================================





% ==========================================================
%                 请在上一行结束填写答案
% ==========================================================

\end{answer}

% ============================================================
%
%              ↑↑↑↑↑↑ 学生答案区域 ↑↑↑↑↑↑
%
% ============================================================


\vspace{1em}



% ============================================================
% ============================================================
%
%                           题目 11
%
% ============================================================
% ============================================================

\begin{problem}[课本习题29]

根据报道美国人血型的分布近似地为：A型为 $37\%$，O型为 $44\%$，B型为 $13\%$，AB型为 $6\%$。夫妻拥有的血型是相互独立的。
\begin{enumerate}[label=(\arabic*)]
    \item B型的人只有输入B、O两种血型才安全。若妻为B型，夫为何种血型未知，求夫是妻的安全输血者的概率。
    \item 随机地取一对夫妇，求妻为B型夫为A型的概率。
    \item 随机地取一对夫妇，求其中一人为A型，另一人为B型的概率。
    \item 随机地取一对夫妇，求其中至少有一人是O型的概率。
\end{enumerate}

\end{problem}


% ============================================================
%
%              ↓↓↓↓↓↓ 学生答案区域 ↓↓↓↓↓↓
%
% ============================================================

\begin{answer}

% ==========================================================
%                 请从下一行开始填写答案
% ==========================================================





% ==========================================================
%                 请在上一行结束填写答案
% ==========================================================

\end{answer}

% ============================================================
%
%              ↑↑↑↑↑↑ 学生答案区域 ↑↑↑↑↑↑
%
% ============================================================


\vspace{1em}



% ============================================================
% ============================================================
%
%                           题目 12
%
% ============================================================
% ============================================================

\begin{problem}[课本习题38]

袋中装有 $m$ 枚正品硬币、$n$ 枚次品硬币（次品硬币的两面均印有国徽），在袋中任取一枚，将它投掷 $r$ 次，已知每次都得到国徽。问这枚硬币是正品的概率为多少？

\end{problem}


% ============================================================
%
%              ↓↓↓↓↓↓ 学生答案区域 ↓↓↓↓↓↓
%
% ============================================================

\begin{answer}

% ==========================================================
%                 请从下一行开始填写答案
% ==========================================================





% ==========================================================
%                 请在上一行结束填写答案
% ==========================================================

\end{answer}

% ============================================================
%
%              ↑↑↑↑↑↑ 学生答案区域 ↑↑↑↑↑↑
%
% ============================================================


\vspace{1em}



% ============================================================
% ============================================================
%
%                           题目 13
%
% ============================================================
% ============================================================

\begin{problem}[课本习题40]

将A、B、C三个字母之一输入信道，输出为原字母的概率为 $\alpha$，而输出为其他一字母的概率都是 $(1-\alpha)/2$。今将字母串 $\mathrm{AAAA}$，$\mathrm{BBBB}$，$\mathrm{CCCC}$ 之一输入信道，输入 $\mathrm{AAAA}$，$\mathrm{BBBB}$，$\mathrm{CCCC}$ 的概率分别为 $p_1,p_2,p_3$（$p_1+p_2+p_3=1$），已知输出为 $\mathrm{ABCA}$，问输入的是 $\mathrm{AAAA}$ 的概率是多少？（设信道传输各个字母的工作是相互独立的。）

\end{problem}


% ============================================================
%
%              ↓↓↓↓↓↓ 学生答案区域 ↓↓↓↓↓↓
%
% ============================================================

\begin{answer}

% ==========================================================
%                 请从下一行开始填写答案
% ==========================================================





% ==========================================================
%                 请在上一行结束填写答案
% ==========================================================

\end{answer}

% ============================================================
%
%              ↑↑↑↑↑↑ 学生答案区域 ↑↑↑↑↑↑
%
% ============================================================



\end{document}
